\addcontentsline{toc}{subsubsection}{Solution 63. The Hidden Euclidean Length}

\begin{dlsummary}[title={\bfseries Solution 63. The Hidden Euclidean Length},breakable]

\begin{proof}
The unique solution is
\[
\boxed{F(a,b)=\sqrt{a^2+b^2}\qquad(a,b\in\mathbb Q).}
\]

\paragraph{1. Recover the basic symmetries.}

Using (2),
\[
F(-1,0)^2=F(1,0)=1,
\]
so \(F(-1,0)=1\). Consequently,
\[
F(-a,-b)=F(a,b).
\]

Similarly,
\[
F(0,1)^2=F(-1,0)=1,
\]
so \(F(0,1)=1\). Thus
\[
F(0,t)=F(t,0)\qquad(t\in\mathbb Q). \tag{3}
\]

\paragraph{2. Determine the values on the axes.}

Repeated use of (1) gives
\[
F(n,0)\le n\qquad(n\in\mathbb Z_{>0}).
\]
Meanwhile, multiplicativity and the normalisation give
\[
F(2^k,0)=2^k.
\]

Fix a positive integer \(n\), and choose \(2^k>n\). By (1),
\[
2^k
\le F(n,0)+F(2^k-n,0)
\le F(n,0)+2^k-n.
\]
Therefore \(F(n,0)\ge n\), proving equality.

For positive integers \(m,n\), multiplicativity now gives
\[
m=F(m,0)=F(m/n,0)F(n,0)=nF(m/n,0).
\]
Using symmetry and (3), we obtain
\[
\boxed{F(t,0)=F(0,t)=|t|\qquad(t\in\mathbb Q).} \tag{4}
\]

In particular, (1) yields the preliminary bound
\[
\boxed{F(a,b)\le |a|+|b|.} \tag{5}
\]

\paragraph{3. Improve the bound by taking powers.}

Define multiplication of pairs by
\[
(a,b)\ast(c,d)=(ac-bd,ad+bc).
\]
For a fixed pair \((a,b)\), write its \(n\)-th power as \((A_n,B_n)\). All coordinates remain rational.

Condition (2) gives
\[
F(A_n,B_n)=F(a,b)^n.
\]
The elementary identity
\[
(ac-bd)^2+(ad+bc)^2=(a^2+b^2)(c^2+d^2)
\]
also gives
\[
A_n^2+B_n^2=(a^2+b^2)^n.
\]

Apply (5) and \((|u|+|v|)^2\le2(u^2+v^2)\):
\[
\begin{aligned}
F(a,b)^n
&\le |A_n|+|B_n|\\
&\le\sqrt2\,(a^2+b^2)^{n/2}.
\end{aligned}
\]
For a nonzero pair, take \(n\)-th roots and let \(n\to\infty\):
\[
\boxed{F(a,b)\le\sqrt{a^2+b^2}.} \tag{6}
\]

\paragraph{4. Use the conjugate pair to force equality.}

Since
\[
(a,b)\ast(a,-b)=(a^2+b^2,0),
\]
equations (2) and (4) give
\[
F(a,b)F(a,-b)=a^2+b^2.
\]
But (6) bounds both factors above by \(\sqrt{a^2+b^2}\). Their product reaches the product of these upper bounds, so both must attain them:
\[
F(a,b)=\sqrt{a^2+b^2}.
\]
The zero pair is already determined.

\paragraph{5. Verify.}

The displayed algebraic identity proves (2). For (1), use
\[
ac+bd\le\sqrt{a^2+b^2}\sqrt{c^2+d^2},
\]
which follows from \((ad-bc)^2\ge0\). Expanding the squares then proves the required triangle inequality.

All normalisations hold.
\end{proof}

\paragraph{Takeaways.}
\begin{itemize}
\item **An exact value can propagate backwards:** Knowing the values at powers of \(2\) determines every integer on an axis.
\item **High powers sharpen estimates:** The crude factor \(\sqrt2\) disappears after taking \(n\)-th roots.
\item **Conjugates produce a known product:** Pairing \((a,b)\) with \((a,-b)\) forces equality in both upper bounds.
\item **Geometry can emerge from algebra:** Neither distance nor length appears in the hypotheses.
\end{itemize}

\end{dlsummary}
